% §12.3 Approximate Joins in Databases
%
% Monograph register: MinHash + LSH similarity join read through the framework.
% Worked-example fix: (b,r)=(20,10) gives J* = (1/20)^(1/10) ~ 0.74, ABOVE the
% threshold theta = 0.5, so it is an aggressive precision-favoring tuning; the
% draft's "most J=0.7 survive" / "half of J=0.6 survive" contradicted the
% (correct) survival numbers 0.434 / ~0.11. Corrected the prose; numbers kept.
% rem:bloom-equijoin salvages ex:ch12-10. Cross-refs sec:minhash-lsh, def:minhash,
% prop:minhash-collision, prop:minhash-error, sec:relation-composition,
% thm:join-rates, def:bloom-filter, ch:approximate-sets, sec:ch12-bridge. DEFINED:
% sec:approximate-joins, rem:bloom-equijoin.

\section{Approximate Joins in Databases}
\label{sec:approximate-joins}

\Cref{sec:relation-composition} treated the join algebraically. This section
reads the database \emph{similarity join} through the same framework, showing the
published MinHash-LSH analyses to be the chapter-12 join rates specialized to a
sketch-based predicate.

\subsection*{The similarity join}

For relations $R_1, R_2$ whose rows are token sets (the words of a document, the
shingles of a string), the similarity join at threshold $\theta$ is $R_1
\bowtie_\theta R_2 = \{(r_1, r_2) : J(r_1, r_2) \ge \theta\}$, the cross-relation
pairs whose Jaccard overlap is at least $\theta$. Record linkage, near-duplicate
detection, and entity resolution all reduce to it. The naive evaluation is $|R_1|
\cdot |R_2|$ Jaccard computations, prohibitive at scale; the framework reads the
join as an approximate relation $\tilde S_\theta(r_1, r_2) = [\hat J(r_1, r_2)
\ge \theta]$ implemented by a sketch.

The MinHash sketch (\Cref{def:minhash}) is exactly that estimator: the fraction
of agreeing signature coordinates is an unbiased estimate of $J$
(\Cref{prop:minhash-collision}). Build the length-$k$ signature for every row
(cost linear in the total token count), and for each candidate pair count
agreements and threshold at $\theta$. The two error rates come from the
chapter-10 Hoeffding bound (\Cref{prop:minhash-error}): a false positive needs
$\hat J$ to overshoot a sub-$\theta$ truth, a false negative needs it to
undershoot a super-$\theta$ truth, and both are bounded by $\Prob{|\hat J - J|
\ge t} \le 2\exp(-2kt^2)$ with $t$ the gap to $\theta$. Larger $k$ shrinks both
exponentially.

\subsection*{LSH bands}

The remaining quadratic factor is candidate enumeration, which LSH
(\Cref{sec:minhash-lsh}) prunes. Partition the signature into $b$ bands of $r$
rows ($k = br$), hash each band, and call $(r_1, r_2)$ a candidate if any band
collides. A band collides with probability $J^r$, so a pair survives with
probability $1 - (1 - J^r)^b$, the S-curve with inflection $J^* \approx
(1/b)^{1/r}$. Placing $J^*$ near $\theta$ concentrates candidates on pairs with
$J \gtrsim \theta$; the candidate set is sub-quadratic when most pairs have $J
\ll \theta$. The pipeline is the intersection of the LSH candidate relation
$\tilde C$ and the threshold relation $\tilde S_\theta$, composing by
\Cref{thm:join-rates}; the LSH layer shrinks the false-positive denominator
(fewer candidates) and adds to the false-negative side (a truly similar pair can
fail to collide), combining by the chapter-3 OR rule.

\begin{remark}[The Bloom equijoin]
\label{rem:bloom-equijoin}
The equijoin, matching rows on equality of a shared key, admits a different
sketch. Build a Bloom filter $B_2$ over $R_2$'s keys, scan $R_1$, and for each
$r_1$ whose key hits $B_2$ scan $R_2$ for the match. The candidate relation ``$r_1$'s
key hits $B_2$'' is a one-sided FNR-zero approximate relation
(\Cref{def:bloom-filter}) with per-pair FPR the Bloom rate; an exact verification
on each candidate then removes the residual false positives, so the \emph{final}
join is exact ($\fprate = \fnrate = 0$) with the Bloom filter serving only as a
candidate pruner. This is the framework's certain-reject structure
(\Cref{ch:approximate-sets}) deployed as a join accelerator, the equijoin
counterpart of the MinHash similarity pipeline.
\end{remark}

\subsection*{Worked example}

Take $|R_1| = |R_2| = 1000$ (so $10^6$ candidate pairs), threshold $\theta =
0.5$, and $k = 200$ partitioned as $b = 20$ bands of $r = 10$. The per-pair
MinHash threshold is sharp: a pair at $J = 0.3$ is emitted with probability
$\le \exp(-2 \cdot 200 \cdot 0.2^2) = \exp(-16) \approx 1.1 \times 10^{-7}$, and a
pair at $J = 0.7$ is missed with the same probability, so away from the boundary
the threshold adds negligible error (at $J = \theta$ exactly it is $1/2$).

The LSH layer dominates, and its tuning needs care. With $(b, r) = (20, 10)$ the
inflection sits at $J^* = (1/20)^{1/10} \approx 0.74$, \emph{above} the threshold
$0.5$, making this an aggressive precision-favoring tuning: the survival
probability $1 - (1 - J^{10})^{20}$ is $\approx 0.434$ at $J = 0.7$, $\approx
0.019$ at the threshold $J = 0.5$, and $\approx 1.2 \times 10^{-4}$ at $J = 0.3$.
Only $43\%$ of strongly similar pairs and $2\%$ of threshold pairs survive, so the
pruning is severe. With a true-similarity distribution of $0.1\%$ at $J \ge 0.5$,
$1\%$ at $[0.3, 0.5)$, and $98.9\%$ below, the candidate set is about
\[
  1000 \cdot 0.434 + 10000 \cdot 0.019 + 989000 \cdot 1.2 \times 10^{-4}
  \approx 434 + 190 + 119 \approx 743
\]
pairs, three orders of magnitude below the $10^6$ all-pairs scan: the
sub-quadratic win. The cost is the low survival of near-threshold pairs, a real
recall penalty; centering $J^*$ on $\theta = 0.5$ (the standard tuning $(1/b)^{1/r}
= \theta$, here a larger $b$) would admit more candidates in exchange for higher
recall. The candidate verification is $O(k)$ per pair, $\approx 1.5 \times 10^5$
signature comparisons against the $2 \times 10^8$ of the naive scan. This
precision/recall trade through $(b, r)$ is the standard operating regime of
MinHash-LSH similarity joins. \Cref{sec:ch12-bridge} closes the chapter.
